\documentclass[12pt,a4paper]{article}
\usepackage{fontspec}
\setmainfont{Times New Roman}
\usepackage[top=2.0cm,bottom=2.0cm,left=2.5cm,right=2.0cm]{geometry}
\usepackage{enumitem}
\usepackage{tabularx}
\usepackage{xltabular}
\usepackage{booktabs}
\usepackage{array}
\newcolumntype{Y}{>{\centering\arraybackslash}X}
\newcolumntype{C}[1]{>{\centering\arraybackslash}p{#1}}
\newcolumntype{L}[1]{>{\raggedright\arraybackslash}p{#1}}
\newcolumntype{R}[1]{>{\raggedleft\arraybackslash}p{#1}}
\pagestyle{plain}
\linespread{1.15}
\setlength{\parskip}{4pt plus 1pt minus 1pt}
\sloppy

\begin{document}
% --- HEADER 2 CỘT CHUẨN NGHỊ ĐỊNH 30/2020/NĐ-CP ---
\noindent
\begin{minipage}[t]{0.46\textwidth}
    \begin{center}
        \fontsize{11pt}{13pt}\selectfont
        CÔNG AN THÀNH PHỐ HÀ NỘI\\
        \textbf{PHÒNG CẢNH SÁT HÌNH SỰ (PC02)}\\[-0.2em]
        \rule{3.2cm}{0.6pt}\\[0.25cm]
        \fontsize{11pt}{13pt}\selectfont Số: \textbf{11/BB-LK}
    \end{center}
\end{minipage}
\hfill
\begin{minipage}[t]{0.52\textwidth}
    \begin{center}
        \fontsize{11pt}{13pt}\selectfont
        \textbf{CỘNG HÒA XÃ HỘI CHỦ NGHĨA VIỆT NAM}\\
        \textbf{Độc lập - Tự do - Hạnh phúc}\\[-0.2em]
        \rule{3.6cm}{0.6pt}\\[0.25cm]
        \textit{Hà Nội, ngày 25 tháng 07 năm 2016}
    \end{center}
\end{minipage}

\vspace{0.5cm}

\begin{center}
    \textbf{\Large BIÊN BẢN LẤY LỜI KHAI}\\[0.15cm]
    \textit{(Lần 2 đối với đối tượng Lê Quang Vũ)}
\end{center}

\vspace{0.3cm}


Vào hồi 20 giờ 30 phút, ngày 25 tháng 07 năm 2016.

Sau khi điều tra viên đưa ra các chứng cứ: Khung ảnh vỡ kính tại hiện trường, Bài báo có nội dung về vụ tai nạn 20 năm trước, Dấu vân tay trên khung ảnh được giám định là của Nguyễn Thanh Tùng. Đối tượng ngay lập tức được triệu tập và tiến hành thẩm vấn lấy lời khai lần hai.


\noindent\textbf{Hỏi:} \textit{Anh Nguyễn Thanh Tùng, anh giải thích như thế nào về dấu vân tay trên khung ảnh tại hiện trường, dù trước đó anh khẳng định không hề gặp mặt nạn nhân? Lý do gì khiến anh hẹn gặp anh Khang và chuyện gì đã diễn ra?}\\[0.15cm]

\noindent\textbf{Đáp:} Thưa cán bộ, tôi xin khai hết sự thật! Suốt 20 năm qua, gia đình tôi chưa bao giờ nguôi ngoai nỗi đau mất đi em trai của mình. Và trớ trêu thay, mới gần đây vì thằng Khang nhỡ mồm lúc say mà tôi mới biết được sự thật: nó chính là người đã dụ em tôi trốn trong tủ và dẫn đến tai nạn thương tâm đó.\\[0.35cm]
Vì lúc đó không được tỉnh táo nên tôi chỉ im lặng ghi nhớ để hẹn nó hôm sau tôi tới nhà để chính tôi đối chất lại với thằng khốn đó.


Chiều tối 24/07 tôi gọi điện hẹn 8h tối qua nhà Khang. Tôi cầm tờ báo qua để chất vấn nó vì sao nhẫn tâm đến mức vứt chiếc còi của Huy ra ngoài tủ như thế. Và tôi muốn nó phải nhận tội, phải cúi đầu cầu xin sự tha thứ của em trai tôi. Nhưng thằng Khang trơ tráo, buông lời cười cợt, xé tờ báo \& cho rằng tất cả chỉ là trò đùa của bọn con nít mà thôi.


Cơn giận bùng lên, tôi lao vào giằng co, xô mạnh thằng Khang ngã đập đầu vào cạnh bàn trà làm đổ vỡ toang bộ bình trà thủy tinh và làm rơi vỡ khung ảnh xuống sàn. Sau đó tôi thấy nó nằm bất động trên sàn nhà. Sợ rằng mình đã gây ra một tai nạn kinh khủng, tôi hoảng loạn bỏ chạy khỏi nhà Khang, chạy một mạch ra đầu ngõ bắt xe ôm về phòng trọ công nhân ở Cầu Bươu.

Nhưng tôi không phải là kẻ đã đâm vào cổ Khang! Tôi cứ tưởng nó chỉ ngất mê man một lúc rồi tự tỉnh lại, ai ngờ sáng hôm sau nghe tin nó bị ai đó đâm chết! Cán bộ ơi, tôi chỉ xô ngã nó thôi chứ tôi thề có vong linh em tôi là tôi không hề giết nó!

\vspace{0.8cm}
\noindent
\begin{minipage}[t]{0.48\textwidth}
    \begin{center}
        \textbf{ĐIỀU TRA VIÊN CHỦ TRÌ}\\[0.1cm]
        \textit{(Ký tên)}\\[1.6cm]
        \textbf{Đại úy Lê Minh}
    \end{center}
\end{minipage}
\hfill
\begin{minipage}[t]{0.48\textwidth}
    \begin{center}
        \textbf{NGƯỜI TỰ THÚ}\\[0.1cm]
        \textit{(Ký, ghi rõ họ tên)}\\[1.6cm]
        \textbf{Nguyễn Thanh Tùng}
    \end{center}
\end{minipage}
\end{document}
